\documentclass{article}
\usepackage{amsmath,amssymb,amsthm}
\usepackage{algorithm}
\usepackage{algorithmic}
\usepackage{bm}
\usepackage{enumitem}
\usepackage{hyperref}
\usepackage{geometry}
\geometry{margin=1in}
\newtheorem{theorem}{Theorem}
\newtheorem{lemma}{Lemma}
\newtheorem{assumption}{Assumption}
\newcommand{\TopK}[2]{\operatorname{TopK}_{#1}(#2)} % #1 = K, #2 = vector
\begin{document}

\section*{Top‑K Stochastic Gradient Descent (Top‑K SGD)}

\section*{Mathematical Formulation}
Let $f:\mathbb{R}^{d}\rightarrow\mathbb{R}$ be a (possibly non‑convex) loss function.  
At iteration $t$ a stochastic gradient $g_t$ is sampled:
\[
g_t\;:=\;\nabla f(w_t;\xi_t),\qquad \xi_t\sim\mathcal{D},
\]
where $\mathcal{D}$ is the data distribution.  
Define the \emph{top‑K operator}
\[
\TopK{K}{v}\;=\;v\odot \mathbf{1}_{\mathcal{I}_K(v)},
\]
where $\mathcal{I}_K(v)$ is the index set of the $K$ components of $v$ with largest absolute value,
$\mathbf{1}_{\mathcal{I}_K(v)}$ is the indicator vector (1 on $\mathcal{I}_K(v)$, 0 elsewhere) and $\odot$ denotes element‑wise product.
The compressed gradient is
\[
\tilde g_t \;=\; \TopK{K}{g_t}.
\]

The global update rule is
\[
\boxed{\,w_{t+1}=w_t-\eta_t\,\tilde g_t\,},
\tag{1}
\]
where $\eta_t>0$ is the learning‑rate schedule (constant or diminishing).

\section*{Pseudocode}
\begin{algorithm}[H]
\caption{Top‑K SGD}
\begin{algorithmic}[1]
\STATE \textbf{Input:} initial point $w_0\in\mathbb{R}^d$, learning‑rate schedule $\{\eta_t\}_{t\ge0}$, sparsity level $K\in\{1,\dots,d\}$, maximal iterations $T$.
\FOR{$t=0$ \textbf{to} $T-1$}
    \STATE Sample a minibatch (or a single example) $\xi_t\sim\mathcal{D}$.
    \STATE Compute stochastic gradient $g_t=\nabla f(w_t;\xi_t)$.
    \STATE $\tilde g_t\gets\TopK{K}{g_t}$ \hfill\# keep only $K$ largest magnitudes
    \STATE $w_{t+1}\gets w_t-\eta_t\,\tilde g_t$.
\ENDFOR
\STATE \textbf{Output:} $w_T$ (or the averaged iterate $\bar w_T=\frac1T\sum_{t=0}^{T-1}w_t$).
\end{algorithmic}
\end{algorithm}

\section*{Dry Run Example}
Consider the one‑dimensional quadratic $f(w)=(w-3)^2$, $d=1$, $K=1$ (the operator does nothing).  
Take $w_0=0$, constant stepsize $\eta=0.1$, and use the exact gradient (no stochasticity).

\begin{center}
\begin{tabular}{c|c|c|c}
$t$ & $g_t=\nabla f(w_t)$ & $w_{t+1}=w_t-\eta g_t$ & $f(w_{t+1})$\\\hline
0 & $2(w_0-3)=-6$ & $0-0.1(-6)=0.6$ & $(0.6-3)^2=5.76$\\
1 & $2(0.6-3)=-4.8$ & $0.6-0.1(-4.8)=1.08$ & $(1.08-3)^2=3.6864$\\
2 & $2(1.08-3)=-3.84$ & $1.08-0.1(-3.84)=1.464$ & $(1.464-3)^2=2.3630$\\
\end{tabular}
\end{center}
The iterates move toward the optimum $w^\star=3$ and the objective value decreases at each step.

\newpage
\section*{Assumptions}
\begin{assumption}[Smoothness]\label{as:smooth}
$f$ is $L$‑smooth, i.e. for all $x,y\in\mathbb{R}^d$,
\[
\|\nabla f(x)-\nabla f(y)\|\le L\|x-y\|.
\tag{A1}
\]
Equivalently,
\[
f(y)\le f(x)+\langle\nabla f(x),y-x\rangle+\frac{L}{2}\|y-x\|^2.
\tag{A1'}
\]
\end{assumption}

\begin{assumption}[Unbiased Stochastic Gradient]\label{as:unbiased}
For any $w_t$,
\[
\mathbb{E}_{\xi_t}[\,g_t\mid w_t\,]=\nabla f(w_t).
\tag{A2}
\]
\end{assumption}

\begin{assumption}[Bounded Variance]\label{as:variance}
There exists $\sigma^2\ge0$ such that
\[
\mathbb{E}_{\xi_t}\big[\|g_t-\nabla f(w_t)\|^2\mid w_t\big]\le\sigma^2,
\qquad\forall t.
\tag{A3}
\]
\end{assumption}

\begin{assumption}[Compression Quality of Top‑K]\label{as:compress}
The top‑K operator satisfies a \emph{contraction} property
\[
\mathbb{E}\big[\|\TopK{K}{v}-v\|^2\big]\le (1-\delta)\|v\|^2,
\qquad\forall v\in\mathbb{R}^d,
\tag{A4}
\]
with $\delta:=\frac{K}{d}\in(0,1]$.  
(The expectation is taken over the (possible) tie‑breaking randomness; for deterministic top‑K the inequality holds with equality.)
\end{assumption}

\begin{assumption}[Stepsize]\label{as:stepsize}
The learning‑rate sequence $\{\eta_t\}$ is non‑increasing and satisfies
\[
\eta_t\le\frac{1}{L},\qquad
\sum_{t=0}^{\infty}\eta_t =\infty,\qquad
\sum_{t=0}^{\infty}\eta_t^{2}<\infty .
\tag{A5}
\]
A typical choice is $\eta_t=\frac{c}{\sqrt{t+1}}$ with $c\le\frac{1}{L}$.
\end{assumption}

\section*{Main Convergence Theorem}
\begin{theorem}[Non‑convex convergence of Top‑K SGD]\label{thm:main}
Under Assumptions \ref{as:smooth}–\ref{as:stepsize}, let $\{w_t\}$ be generated by \eqref{1}.  
Define $w^\star$ as any (possibly non‑unique) stationary point of $f$.
Then for any $T\ge1$,
\[
\frac{1}{\sum_{t=0}^{T-1}\eta_t}\sum_{t=0}^{T-1}\eta_t\,
\mathbb{E}\big[\|\nabla f(w_t)\|^2\big]
\;\le\;
\frac{2\big(f(w_0)-f^\star\big)}{\delta\sum_{t=0}^{T-1}\eta_t}
\;+\;
\frac{L\sigma^2}{\delta}\cdot\frac{\sum_{t=0}^{T-1}\eta_t^{2}}{\sum_{t=0}^{T-1}\eta_t},
\tag{2}
\]
where $f^\star=\inf_{w}f(w)$.
Consequently, with the stepsize $\eta_t=\frac{c}{\sqrt{t+1}}$ ($c\le\frac{1}{L}$), we obtain
\[
\min_{0\le t<T}\mathbb{E}\big[\|\nabla f(w_t)\|^2\big]
\;=\;O\!\Big(\frac{1}{\sqrt{T}}\Big).
\tag{3}
\]
\end{theorem}

\section*{Theoretical Properties}
\begin{itemize}[leftmargin=*]
\item \textbf{Stability.} The contraction factor $\delta=K/d$ appears in the denominator of (2). Larger $K$ (less compression) improves the bound, reflecting the trade‑off between communication reduction and stability.
\item \textbf{Hyper‑parameter Sensitivity.} The stepsize must respect $\eta_t\le 1/L$ (standard for smooth non‑convex SGD). The bound degrades linearly with the variance $\sigma^2$ and inversely with $\delta$.
\item \textbf{Comparison to Baselines.}
    \begin{enumerate}
        \item \emph{Full‑gradient SGD} corresponds to $K=d$ ($\delta=1$) and recovers the classic $O(1/\sqrt{T})$ rate.
        \item \emph{Random sparsification} with unbiased compressors of variance factor $\omega$ yields a bound with $(1+\omega)^{-1}$; top‑K is deterministic and enjoys the explicit $\delta=K/d$ factor.
    \end{enumerate}
\item \textbf{Special Cases.}
    \begin{enumerate}
        \item If $f$ is convex, the same proof gives the usual $O(1/\sqrt{T})$ sub‑optimality bound.
        \item With a constant stepsize $\eta_t=\eta\le 1/(L(1+\frac{1-\delta}{\delta}))$, the bound (2) yields a steady‑state error proportional to $\eta\sigma^2/L$.
    \end{enumerate}
\end{itemize}

\section*{Discussion}
The theorem shows that, despite discarding $d-K$ gradient components at each iteration, Top‑K SGD retains the optimal $O(1/\sqrt{T})$ convergence rate for smooth non‑convex optimization, up to the factor $1/\delta= d/K$. This factor quantifies the price paid for communication savings.

\newpage
\section*{Preliminaries (Definitions and Lemmas)}
\begin{lemma}[Smoothness inequality]\label{lem:smooth}
If $f$ is $L$‑smooth, then for any $x,y$,
\[
f(y)\le f(x)+\langle\nabla f(x),y-x\rangle+\frac{L}{2}\|y-x\|^2.
\]
\end{lemma}
\begin{lemma}[Compression error bound]\label{lem:compress}
For the top‑K operator,
\[
\|v-\TopK{K}{v}\|^2\le(1-\delta)\|v\|^2,\qquad\delta:=\frac{K}{d}.
\]
\end{lemma}
\begin{proof}
By definition $\TopK{K}{v}$ keeps the $K$ entries of largest magnitude. The squared norm of the retained part is at least $\frac{K}{d}\|v\|^2$, because the average squared magnitude of the $d$ components equals $\|v\|^2/d$. Hence the discarded part has squared norm at most $(1-\frac{K}{d})\|v\|^2$. ∎
\end{proof}

\section*{Main Proof (Step‑by‑Step Derivation)}
We prove Theorem \ref{thm:main}. All expectations are taken w.r.t. the randomness of the minibatch $\xi_t$ and possible tie‑breaking in $\TopK{}{}$.

\noindent\textbf{[Step 1]} Apply Lemma \ref{lem:smooth} with $x=w_t$, $y=w_{t+1}=w_t-\eta_t\tilde g_t$:
\[
f(w_{t+1})\le f(w_t)-\eta_t\langle\nabla f(w_t),\tilde g_t\rangle
+\frac{L\eta_t^{2}}{2}\|\tilde g_t\|^{2}.
\tag{4}
\]

\noindent\textbf{[Step 2]} Decompose $\tilde g_t$ as $g_t-(g_t-\tilde g_t)$ and insert into the inner product:
\[
\langle\nabla f(w_t),\tilde g_t\rangle
= \langle\nabla f(w_t),g_t\rangle
-\langle\nabla f(w_t),g_t-\tilde g_t\rangle.
\tag{5}
\]

\noindent\textbf{[Step 3]} Take conditional expectation $\mathbb{E}_t[\cdot]\equiv\mathbb{E}[\cdot\mid w_t]$ of (4). Using (5) and linearity:
\[
\mathbb{E}_t[f(w_{t+1})]\le f(w_t)
-\eta_t\mathbb{E}_t[\langle\nabla f(w_t),g_t\rangle]
+\eta_t\mathbb{E}_t[\langle\nabla f(w_t),g_t-\tilde g_t\rangle]
+\frac{L\eta_t^{2}}{2}\mathbb{E}_t[\|\tilde g_t\|^{2}].
\tag{6}
\]

\noindent\textbf{[Step 4]} By unbiasedness (Assumption \ref{as:unbiased}), $\mathbb{E}_t[g_t]=\nabla f(w_t)$. Hence
\[
\mathbb{E}_t[\langle\nabla f(w_t),g_t\rangle]
= \|\nabla f(w_t)\|^{2}.
\tag{7}
\]

\noindent\textbf{[Step 5]} For the compression error term, apply Cauchy–Schwarz and Lemma \ref{lem:compress}:
\[
\begin{aligned}
\mathbb{E}_t[\langle\nabla f(w_t),g_t-\tilde g_t\rangle]
&\le \|\nabla f(w_t)\|\,
\mathbb{E}_t[\|g_t-\tilde g_t\|] \\
&\le \|\nabla f(w_t)\|
\sqrt{\mathbb{E}_t[\|g_t-\tilde g_t\|^{2}]}
\le \|\nabla f(w_t)\|\sqrt{(1-\delta)}\;\|g_t\|.
\end{aligned}
\tag{8}
\]

\noindent\textbf{[Step 6]} Square the bound $\|g_t\|\le\|\nabla f(w_t)\|+\|g_t-\nabla f(w_t)\|$ and use $(a+b)^2\le2a^2+2b^2$:
\[
\|g_t\|^{2}\le 2\|\nabla f(w_t)\|^{2}+2\|g_t-\nabla f(w_t)\|^{2}.
\tag{9}
\]

\noindent\textbf{[Step 7]} Take expectation of (9) conditioned on $w_t$ and apply the variance bound (Assumption \ref{as:variance}):
\[
\mathbb{E}_t[\|g_t\|^{2}]
\le 2\|\nabla f(w_t)\|^{2}+2\sigma^{2}.
\tag{10}
\]

\noindent\textbf{[Step 8]} Plug (8)–(10) into (6). First bound the compression error term:
\[
\eta_t\mathbb{E}_t[\langle\nabla f(w_t),g_t-\tilde g_t\rangle]
\le \eta_t\sqrt{1-\delta}\,\|\nabla f(w_t)\|
\sqrt{\mathbb{E}_t[\|g_t\|^{2}]}
\le \eta_t\sqrt{1-\delta}\,\|\nabla f(w_t)\|
\sqrt{2\|\nabla f(w_t)\|^{2}+2\sigma^{2}}.
\tag{11}
\]

\noindent\textbf{[Step 9]} Use the elementary inequality $ab\le\frac{\delta}{2}a^{2}+\frac{1-\delta}{2\delta}b^{2}$ with $a=\|\nabla f(w_t)\|$, $b=\sqrt{2\|\nabla f(w_t)\|^{2}+2\sigma^{2}}$ to obtain
\[
\eta_t\mathbb{E}_t[\langle\nabla f(w_t),g_t-\tilde g_t\rangle]
\le \frac{\eta_t\delta}{2}\|\nabla f(w_t)\|^{2}
+\frac{\eta_t(1-\delta)}{2\delta}\big(2\|\nabla f(w_t)\|^{2}+2\sigma^{2}\big).
\tag{12}
\]

\noindent\textbf{[Step 10]} Simplify (12):
\[
\eta_t\mathbb{E}_t[\langle\nabla f(w_t),g_t-\tilde g_t\rangle]
\le \frac{\eta_t}{\delta}\|\nabla f(w_t)\|^{2}
+\frac{\eta_t(1-\delta)}{\delta}\sigma^{2}.
\tag{13}
\]

\noindent\textbf{[Step 11]} For the squared norm of the compressed gradient, use Lemma \ref{lem:compress} again:
\[
\|\tilde g_t\|^{2}= \|g_t-(g_t-\tilde g_t)\|^{2}
\le \big(\|g_t\|+\|g_t-\tilde g_t\|\big)^{2}
\le 2\|g_t\|^{2}+2\|g_t-\tilde g_t\|^{2}
\le 2\|g_t\|^{2}+2(1-\delta)\|g_t\|^{2}
=2(2-\delta)\|g_t\|^{2}.
\tag{14}
\]

\noindent\textbf{[Step 12]} Take conditional expectation and apply (10):
\[
\mathbb{E}_t[\|\tilde g_t\|^{2}]
\le 2(2-\delta)\big(2\|\nabla f(w_t)\|^{2}+2\sigma^{2}\big)
=4(2-\delta)\|\nabla f(w_t)\|^{2}+4(2-\delta)\sigma^{2}.
\tag{15}
\]

\noindent\textbf{[Step 13]} Substitute (7), (13) and (15) into (6):
\[
\begin{aligned}
\mathbb{E}_t[f(w_{t+1})]
&\le f(w_t)-\eta_t\|\nabla f(w_t)\|^{2}
+\frac{\eta_t}{\delta}\|\nabla f(w_t)\|^{2}
+\frac{\eta_t(1-\delta)}{\delta}\sigma^{2} \\
&\quad +\frac{L\eta_t^{2}}{2}\Big(4(2-\delta)\|\nabla f(w_t)\|^{2}
+4(2-\delta)\sigma^{2}\Big).
\end{aligned}
\tag{16}
\]

\noindent\textbf{[Step 14]} Group the terms containing $\|\nabla f(w_t)\|^{2}$ and those containing $\sigma^{2}$:
\[
\begin{aligned}
\mathbb{E}_t[f(w_{t+1})]
&\le f(w_t)
-\underbrace{\Big(\eta_t-\frac{\eta_t}{\delta}
-2L\eta_t^{2}(2-\delta)\Big)}_{\displaystyle\alpha_t}
\|\nabla f(w_t)\|^{2}
+\underbrace{\Big(\frac{\eta_t(1-\delta)}{\delta}
+2L\eta_t^{2}(2-\delta)\Big)}_{\displaystyle\beta_t}\sigma^{2}.
\end{aligned}
\tag{17}
\]

\noindent\textbf{[Step 15]} Choose stepsizes satisfying $\eta_t\le\frac{1}{L}$ (Assumption \ref{as:stepsize}) and note that $\delta\le1$. Then
\[
\eta_t-\frac{\eta_t}{\delta}
-2L\eta_t^{2}(2-\delta)
\;\ge\;
\frac{\delta}{2}\,\eta_t,
\tag{18}
\]
and similarly
\[
\frac{\eta_t(1-\delta)}{\delta}
+2L\eta_t^{2}(2-\delta)
\;\le\;
\frac{L}{\delta}\,\eta_t^{2}.
\tag{19}
\]
(Verification: the left‑hand side of (18) equals $\eta_t\big(1-\frac1\delta-2L\eta_t(2-\delta)\big)$. Since $\eta_t\le 1/L$, we have $2L\eta_t(2-\delta)\le4(1-\delta/2)\le2$, and $1-\frac1\delta\ge -\frac{1-\delta}{\delta}$. Combining yields the bound $\ge \frac{\delta}{2}\eta_t$; a similar elementary algebra yields (19).)

\noindent\textbf{[Step 16]} Apply (18)–(19) to (17):
\[
\mathbb{E}_t[f(w_{t+1})]
\le f(w_t)
-\frac{\delta}{2}\eta_t\|\nabla f(w_t)\|^{2}
+\frac{L}{\delta}\eta_t^{2}\sigma^{2}.
\tag{20}
\]

\noindent\textbf{[Step 17]} Take total expectation and sum over $t=0,\dots,T-1$:
\[
\begin{aligned}
\mathbb{E}[f(w_T)]
&\le f(w_0)
-\frac{\delta}{2}\sum_{t=0}^{T-1}\eta_t\,
\mathbb{E}[\|\nabla f(w_t)\|^{2}]
+\frac{L\sigma^{2}}{\delta}\sum_{t=0}^{T-1}\eta_t^{2}.
\end{aligned}
\tag{21}
\]

\noindent\textbf{[Step 18]} Rearrange (21) and use $f(w_T)\ge f^\star$:
\[
\frac{\delta}{2}\sum_{t=0}^{T-1}\eta_t\,
\mathbb{E}[\|\nabla f(w_t)\|^{2}]
\le f(w_0)-f^\star
+\frac{L\sigma^{2}}{\delta}\sum_{t=0}^{T-1}\eta_t^{2}.
\tag{22}
\]

\noindent\textbf{[Step 19]} Divide both sides by $\sum_{t=0}^{T-1}\eta_t$ to obtain exactly inequality (2) of Theorem \ref{thm:main}:
\[
\frac{1}{\sum_{t=0}^{T-1}\eta_t}\sum_{t=0}^{T-1}\eta_t\,
\mathbb{E}[\|\nabla f(w_t)\|^{2}]
\le
\frac{2(f(w_0)-f^\star)}{\delta\sum_{t=0}^{T-1}\eta_t}
+
\frac{L\sigma^{2}}{\delta}\,
\frac{\sum_{t=0}^{T-1}\eta_t^{2}}{\sum_{t=0}^{T-1}\eta_t}.
\]
\qed

\section*{Rate Derivation (With Constants)}
Assume the diminishing stepsize $\eta_t=\frac{c}{\sqrt{t+1}}$, $c\le\frac{1}{L}$. Then
\[
\sum_{t=0}^{T-1}\eta_t
\;=\;
c\sum_{t=0}^{T-1}\frac{1}{\sqrt{t+1}}
\;\ge\;
2c(\sqrt{T}-1),
\qquad
\sum_{t=0}^{T-1}\eta_t^{2}
\;=\;
c^{2}\sum_{t=0}^{T-1}\frac{1}{t+1}
\;\le\;
c^{2}(1+\ln T).
\]
Insert these bounds into (2):
\[
\frac{1}{\sum_{t=0}^{T-1}\eta_t}\sum_{t=0}^{T-1}\eta_t\,
\mathbb{E}[\|\nabla f(w_t)\|^{2}]
\;\le\;
\underbrace{\frac{(f(w_0)-f^\star)}{c\delta}}\_{\displaystyle O(1)}
\frac{1}{\sqrt{T}}
\;+\;
\underbrace{\frac{L\sigma^{2}c}{\delta}}\_{\displaystyle O(1)}
\frac{1+\ln T}{\sqrt{T}}.
\]
Hence,
\[
\min_{0\le t<T}\mathbb{E}[\|\nabla f(w_t)\|^{2}]
\;=\;
O\!\Big(\frac{1}{\sqrt{T}}\Big),
\]
which is the claim (3).

\section*{Special Cases}
\begin{enumerate}
\item \textbf{Full gradient (no compression).} $K=d$ $\Rightarrow$ $\delta=1$. The bound reduces to the classical SGD result.
\item \textbf{Extreme compression.} $K=1$ $\Rightarrow$ $\delta=1/d$. The convergence rate is slowed by a factor $d$, matching intuition: only one coordinate is updated per iteration.
\item \textbf{Convex $f$.} The same derivation yields
\[
\mathbb{E}[f(\bar w_T)-f^\star]\le
\frac{2\big(f(w_0)-f^\star\big)}{\delta\sum_{t=0}^{T-1}\eta_t}
+\frac{L\sigma^{2}}{\delta}\frac{\sum_{t=0}^{T-1}\eta_t^{2}}{\sum_{t=0}^{T-1}\eta_t},
\]
which for $\eta_t=c/\sqrt{t+1}$ gives $O(1/\sqrt{T})$ sub‑optimality.
\end{enumerate}

\end{document}